\section{Erd\H{o}s Problem \#230: a gap above the Parseval bound}
\noindent Partial reconstruction and a finite-degree improvement, 23 September 2026.

\subsection*{1) Statement and known resolution}
Let $P(z)=\sum_{k=1}^n a_kz^k$, where $|a_k|=1$ and $n\ge2$. Does there exist an absolute $c>0$ such that every such polynomial satisfies
\[
 M:=\max_{|z|=1}|P(z)|\ge(1+c)\sqrt n?\tag{1}
\]
Kahane's ultraflat polynomials give a negative answer. Here we prove the stronger finite-degree lower bound $M\ge\sqrt{n+2}$ and a concentration estimate. They are consistent with the known negative answer, because their relative gap tends to zero.

When stating ultraflatness, the quantity close to $\sqrt n$ is $|P(z)|$, not the complex number $P(z)$. Uniform approximation of $P(z)$ itself to a positive constant on the whole circle would contradict its zero constant coefficient, by integration.

\subsection*{2) Isolating the outermost correlation}
The trigonometric polynomial $|P(e^{i\theta})|^2$ has constant coefficient $n$ and frequencies between $-(n-1)$ and $n-1$. Its coefficient at the top frequency $d=n-1$ is
\[
 c_d=a_n\overline{a_1},\qquad |c_d|=1.
\]
Average over a rotated regular $d$-gon. The root-of-unity identity kills every frequency not divisible by $d$, leaving only $0,d,-d$:
\[
 \frac1d\sum_{j=0}^{d-1}
 \left|P\left(e^{i(\theta+2\pi j/d)}\right)\right|^2
 =n+2\operatorname{Re}(c_de^{id\theta}).\tag{2}
\]
This is valid also at $d=1$, when the average has one term. Choose $\theta$ with $c_de^{id\theta}=1$. At least one of the values averaged in (2) is at least $n+2$, proving
\[
 \boxed{\ M\ge\sqrt{n+2}\ }.\tag{3}
\]
For $n=2$ the exact maximum is $2=\sqrt{n+2}$. At $n=3$, the polynomial $P(z)=z+iz^2+z^3$ has
\[
 |P(e^{i\theta})|^2=|2\cos\theta+i|^2=4\cos^2\theta+1,
\]
so its maximum is $\sqrt5$, again equality in (3).

\subsection*{3) What an almost minimal maximum would imply}
Normalize Lebesgue measure on the circle to total mass one. Parseval gives $\int|P|^2=n$. Suppose $M\le(1+\delta)\sqrt n$, and let $E$ be the set where $|P|<(1-\varepsilon)\sqrt n$, with $0<\varepsilon<1$ and $\delta\ge0$. Writing $p=|E|$, the integral bound gives
\[
 1\le p(1-\varepsilon)^2+(1-p)(1+\delta)^2,
\]
and therefore
\[
 p\le
 \frac{(1+\delta)^2-1}{(1+\delta)^2-(1-\varepsilon)^2}.\tag{4}
\]
Thus if the maximum approaches the Parseval scale, almost all circle values approach that scale from below for each fixed $\varepsilon$. This does not guarantee a uniform pointwise lower bound, and it provides no phases $a_k$ achieving such a maximum.

\subsection*{4) The precise gap to the known counterexample}
Bound (3) supplies only the degree-dependent relative improvement $\sqrt{1+2/n}-1\to0$. It cannot establish (1) with a fixed $c$. Conversely, to disprove (1) constructively one needs unimodular coefficient lists of unbounded lengths with $M/\sqrt n\to1$. The stronger known ultraflat construction makes $\bigl||P(z)|/\sqrt n-1\bigr|\to0$ uniformly on the circle, which would immediately suffice, but its coefficient construction and uniform estimates have not been reconstructed here.

\textbf{UNRESOLVED in this attempt.} The finite-degree improvement, its sharp small cases, and (4) are proved; the essential asymptotic counterexample is missing.

\subsection*{5) Sources}
\url{https://www.erdosproblems.com/230}, checked 23 September 2026; E. Bombieri and J. Bourgain, \emph{On Kahane's ultraflat polynomials}, J. Eur. Math. Soc. 11 (2009), 627--703, \url{https://doi.org/10.4171/JEMS/163}. The latter is a primary account of the known construction and improvements, not a construction supplied by this attempt.

The estimate describes progress toward a counterexample, not the published problem status or a correctness probability.
\subsection*{6) COMPLETION ESTIMATE}
\textbf{COMPLETION: 35\%}
